\chapter{Conclusiones}
En el presente trabajo se desarrolló el cálculo cuántico completo del oscilador armónico, obteniendo explícitamente su espectro energético y sus autofunciones, y analizando su significado físico dentro de la mecánica cuántica.\\

\textbf{Espectro de energía cuantizado:} Se derivó analíticamente el espectro discreto de energía $E_n = \hbar\omega(n + \frac{1}{2})$ mediante el formalismo algebraico de operadores escalera de \cite{dirac1930}. La cuantización surge naturalmente de las relaciones de conmutación $[\hat{a}, \hat{a}^\dagger] = 1$ y la condición de no negatividad del operador número $\langle n|\hat{N}|n\rangle = ||\hat{a}|n\rangle||^2 \geq 0$, estableciendo que los niveles energéticos están igualmente espaciados por $\hbar\omega$ y existe una energía de punto cero de $\frac{1   }{2}\hbar\omega$ incluso en el estado fundamental.\\
    
    
\textbf{Autofunciones y propiedades:} Se obtuvieron las autofunciones explícitas $\psi_n(x)$ expresadas en términos de polinomios de Hermite multiplicados por una envolvente gaussiana. Se caracterizaron sus propiedades fundamentales: ortonormalidad (asegurando que los estados son físicamente distinguibles), paridad bien definida según $(-1)^n$, y estructura nodal donde el $n$-ésimo estado presenta exactamente $n$ ceros, confirmando la relación entre nivel de excitación y complejidad espacial de la función de onda.\\

\textbf{Equivalencia de formalismos:} Se demostró la equivalencia matemática entre el método algebraico de \cite{dirac1930} y el método analítico de \cite{schrodinger1926a}, validando que ambos conducen a las mismas autofunciones para el estado fundamental y estados excitados. Esta equivalencia confirma la consistencia interna de la mecánica cuántica y evidencia que diferentes representaciones matemáticas describen la misma física subyacente. El método algebraico resulta más elegante y generalizable, mientras que el analítico proporciona visualización directa de las densidades de probabilidad.\\

\textbf{Importancia del oscilador armónico:} El sistema estudiado trasciende su valor como ejercicio académico. El formalismo de operadores de creación y aniquilación desarrollado constituye el lenguaje natural para describir excitaciones cuantizadas (fonones, fotones, plasmones) en sistemas de muchos cuerpos, y es fundamental en aplicaciones contemporáneas como circuit QED y códigos bosónicos para corrección de errores en computación cuántica. La solución del oscilador armónico cuántico representa así un puente conceptual entre los fundamentos teóricos establecidos por \cite{schrodinger1926a} y \cite{dirac1930} en la década de 1920 y las tecnologías cuánticas del siglo XXI.
